\section{The actual even corner and the rational rank-two frame}
\label{sec:cubic-endpoint-transfer}

The odd-entry ideal detects the endpoint block, but it does not
make the middle character trivial.  This section keeps that character
as a separate formal factor, proves descent of the endpoint block to
the original rational-character inertia quotient, and transfers the
literal rank-two finite-support assertion to every even prime.
The arithmetic hypothesis in that last transfer is named explicitly.
It is not inferred from a determinant-ring map or from the existence
of a cubic seed.

Retain the actual host \(H/F\), \(p=269\), and residual data of
AC-G93--AC-G99.  On the absolute groups put
\[
 \theta=(\chi^j\gamma,1,\chi^j\gamma^{-1}),\quad j=1,2,
 \qquad \chi^3=1,\quad \chi^c=\chi^{-1},\quad
 \gamma^c=\gamma,\quad \gamma(c)=1.
 \tag{G106.1}
\]
Here \(\gamma=\bar\psi\bar\epsilon^e\) is the actual rational
residual character.  Write \(u=[\chi^j]\) for its indicated cubic
Teichmuller lift, so that \(u(g^\star)=u(g)\).
Let \(N_\theta\) and \(N_\gamma\) be the closed normal subgroups
of \(G_F\) generated, respectively, by
\[
 I_w\cap\ker\chi\cap\ker\gamma,
       \qquad I_w\cap\ker\gamma
       \quad(w\nmid p\text{ a place of }H).
 \tag{G106.2}
\]
The first is exactly the joint residual kernel in AC-G97:
if \(\chi^j\gamma=\chi^j\gamma^{-1}=1\), then
\(\gamma^2=1\) and \(\chi^j=\gamma^{-1}\), so both values are one.
Set
\[
 G_\theta=G_H/N_\theta,
 \quad G_\gamma=G_H/N_\gamma,
 \quad \pi_\gamma:G_\theta\longtwoheadrightarrow G_\gamma.
 \tag{G106.3}
\]
All conjugate inertia groups are included in (G106.2).
Both quotients therefore carry the same star operation, cyclotomic
character and full coefficient-prime local groups.  These full
local groups, rather than their images, are used throughout.

Let \(R_\theta,B_\theta,A_\theta\) be the actual tuple rings,
and let \(E_R\subset R_\theta\), \(E_A\subset A_\theta\) be the
ideals generated by all entries of odd weight.  These are exactly
the four endpoint--middle matrix positions.  Put
\(E_B=E_A\cap B_\theta\).
Let \(R_2,B_2,A_2\) denote the same normalized ordinary construction
for \(G_\gamma\) and residual tuple \((\gamma,\gamma^{-1})\),
with weights \((1,-1)\).  It is the literal old rank-two functor:
its residual diagonal is the two powers of the single character
\(\gamma\).  Its existence follows either from the old
\texttt{Space} theorem once the frame conditions below are imposed,
or directly from the general universal-lift construction and the
rank-two version of the Hensel flag argument in AC-G94.
The local residual characters are distinct at the retained inertia
witness, since \(268\nmid2e\).
The same arguments as AC-G95 give its complete weight-zero ring and
finite graded algebraization.  This application makes no change to
its residual condition.

\begin{lemma}[Finite prime-to-\(p\) scalar rigidity]
\label{lem:G106scalar}
Let \(C\) be a local \(\mathcal O\)-algebra, \(d\) prime to \(p\),
and \(T\in\operatorname{GL}_r(C)\) satisfy \(T^d=1\).
If \(T\bmod\mathfrak m_C=\alpha I\), then
\(T=[\alpha]I\).  The assertion allows nilpotents and
\(p\)-torsion in \(C\).
\end{lemma}
\begin{proof}
Factor
\(X^d-1=(X-[\alpha])Q(X)\) over \(\mathcal O\).
The reduction of \(Q(T)\) is the scalar
\(d\alpha^{d-1}\), which is nonzero.  Hence its determinant
is a unit, and \(Q(T)\) is invertible by the adjugate identity.
The polynomial identity
\((T-[\alpha]I)Q(T)=0\) proves the assertion.
\end{proof}

\begin{proposition}[AC-G106-corner: the exact even formal factor and actual algebra map]
\label{AC-G106-corner}
Let \(\mathcal M^-\) be the maximal abelian pro-\(p\) quotient
of \(G_H\) which is unramified at every finite place away from
\(p\), with the additional relations \(cgc^{-1}=g^{-1}\).
It is a finitely generated \(\mathbb Z_p\)-module.  Its completed
group algebra \(R_{\rm mid}=\mathcal O[[\mathcal M^-]]\) represents
continuous characters \(m\equiv1\) satisfying
\(m(g^\star)=m(g)\) and unramifiedness away from \(p\).
There is an exact isomorphism of complete local \(\mathcal O\)-algebras
\[
 R_\theta/E_R
       \simeq R_2\widehat\otimes_{\mathcal O}R_{\rm mid}.
 \tag{G106.4}
\]
On representations it is the inverse pair of constructions
\[
 r=\begin{pmatrix}B_{00}&0&B_{02}\\0&m&0\\B_{20}&0&B_{22}\end{pmatrix}
       \longmapsto (s=u^{-1}B,m),
 \qquad
 (s,m)\longmapsto u s\oplus m
 \tag{G106.5}
\]
where the second direct sum is placed in the ordered positions
\(0,2\) and \(1\).
The representation \(s\) factors through the actual quotient
\(G_\gamma\), has the original rank-two polarization and
normalization, and has local characters
\(\lambda_{v,0}=u^{-1}\kappa_{v,0}\),
\(\lambda_{v,1}=u^{-1}\kappa_{v,2}\).
The middle character on every local group is \(\kappa_{v,1}\).

The factor map restricts to actual algebra homomorphisms
\[
 B_2\longrightarrow \overline B_\theta:=B_\theta/E_B,
 \qquad
 A_2\longrightarrow \overline A_\theta:=A_\theta/E_A,
 \tag{G106.6}
\]
compatible with their inclusions, with the displayed entry and local
character formulas.  The ring \(\overline B_\theta\) is the closed
subring generated by the image of \(B_2\) and the values of \(m\).
No passage to a reduced ring or characteristic-zero fiber is used.
\end{proposition}

\begin{proof}
\emph{1. Block decomposition and the actual inertia quotient.}
Over every complete local test ring, killing the odd entries makes
matrix multiplication block diagonal in the two indicated blocks.
Thus \(B\) is a rank-two representation and \(m\) is a character.
Polarity gives
\[
 B(g^\star)=J_2B(g)^tJ_2,
 \qquad m(g^\star)=m(g).
 \tag{G106.7}
\]
At every away-\(p\) inertia group, its image in \(G_\theta\)
factors through the finite image of \((\chi,\gamma)\), of order
prime to \(p\).  If \(h\) is in that inertia and \(\gamma(h)=1\),
then the residual endpoint block is \(\chi^j(h)I\).
Lemma~\ref{lem:G106scalar} shows that the actual block is
\(u(h)I\).  Hence \(u^{-1}B\) kills every generator of
\(N_\gamma/N_\theta\).  Its continuous kernel is closed and
normal, so it kills their closure and factors through (G106.3).
The same lemma applied to the residual-trivial middle character
shows that \(m\) kills all away-\(p\) inertia.
This proves actual kernel descent, not only equality of semisimple
inertia characters.

Since \(u(g^\star)=u(g)\), untwisting preserves (G106.7).
It also preserves the diagonal normalization at \(g_0\): the
scalar at \(g_0\) and \(g_0^\star\) is the same.
The endpoint weights of rank three are literally the old rank-two
weights \((1,-1)\).

\emph{2. Remove the middle factor from the ordered local identity.}
At every coefficient-prime local group choose a tame inertia element
\(h_*\) on which \(\gamma=1\) and \(\chi\) has order three.
Such an element exists in the actual local fields: the cubic
character has order three, the image of \(\gamma\) has order dividing
\(268\), and taking a suitable prime-to-three power kills the latter
while preserving the former.  Restriction through the retained
\(p\)-power local extensions preserves these finite images.
Consequently
\[
 m(h_*)-\kappa_{v,0}(h_*),\quad
 m(h_*)-\kappa_{v,2}(h_*),\quad
 \det\bigl(B(h_*)-\kappa_{v,1}(h_*)I\bigr)
 \tag{G106.8}
\]
are units in every local test ring.
The middle entry of the three-factor product, evaluated at
\((h_*,h,h_*)\), gives
\(m(h)=\kappa_{v,1}(h)\) for every \(h\).

For completeness the noncommutative two-by-two cancellation is as
follows.  Set \(P=B(x)-\kappa_{v,0}(x)I\),
\(M=B(h_*)-\kappa_{v,1}(h_*)I\), and
\(Q_y=B(y)-\kappa_{v,2}(y)I\).
The ordered product gives \(PMQ_y=0\) for every \(y\).
Thus \(PMB(y)=\kappa_{v,2}(y)PM\), and
\(PM^2=(\kappa_{v,2}(h_*)-\kappa_{v,1}(h_*))PM\).
The two-by-two identity
\(M(\operatorname{tr}(M)I-M)=\det(M)I\) implies
\[
 \det(M)P=
 \bigl(\operatorname{tr}(M)-\kappa_{v,2}(h_*)+
                  \kappa_{v,1}(h_*)\bigr)PM.
 \tag{G106.9}
\]
Multiply by \(Q_y\) and invert \(\det(M)\).  This proves
\(PQ_y=0\), exactly the old ordered rank-two condition.
Scalar untwisting gives the condition for \(s\) and \(\lambda\).
This is the complete elementary argument in the pinned
\path{FrobeniusCornerOrderSkips.lean}; no commutation of \(M\)
with \(P\) has been assumed.

Conversely, if \((s(x)-\lambda_0(x)I)
(s(z)-\lambda_1(z)I)=0\) for every \(x,z\), then
\((s(x)-\lambda_0(x)I)s(y)=\lambda_1(y)
(s(x)-\lambda_0(x)I)\).  Inserting any scalar middle factor
between these two ordered factors still gives zero.  Multiplying
by the three scalar values of \(u\) proves the endpoint part of
the rank-three identity for \(u s\oplus m\); its middle part has
a zero middle factor.  All residual and normalized conditions are
preserved.  Thus (G106.5) gives inverse natural transformations.

\emph{3. Represent the remaining character.}
The quotient \(\mathcal M^-\) is a quotient of \(G_\theta\).
The arithmetic \(\Phi_p\) condition of AC-G96 makes
\(\operatorname{Hom}_{\rm cont}(\mathcal M^-,\mathbb F_p)\)
finite dimensional.  Equivalently, \(\mathcal M^-/p\mathcal M^-\)
is finite dimensional; choose elements of \(\mathcal M^-\) lifting
a basis of this quotient.
The resulting finitely generated closed subgroup has quotient
\(C\) with \(C=pC\).  A nonzero compact abelian pro-\(p\)
group has a nonzero finite \(p\)-quotient and then a nonzero
\(\mathbb F_p\)-quotient, contradicting \(C=pC\).
Hence \(\mathcal M^-\) is finitely generated over \(\mathbb Z_p\).
Its structure is \(\mathbb Z_p^r\oplus T\) with \(T\) finite;
\(\mathcal O[[\mathcal M^-]]\) is finite over
\(\mathcal O[[X_1,\ldots,X_r]]\) and is local, since the
augmentation ideal of \(k[T]\) is nilpotent.
Its usual continuous group-algebra property is precisely the
character functor in the statement.  All its target characters
have pro-\(p\) image, because \(1+\mathfrak m_C\) is the inverse
limit of finite \(p\)-groups.
The completed tensor product represents the product of the two
complete-local functors.  The natural transformations just proved
and their inverse establish (G106.4).

\emph{4. The algebraic quotient injects into its completion.}
We supply this point because a map between complete rings alone
would not give the maps (G106.6) at arbitrary primes of \(A_\theta\).
Let \(x_{\rm diag}\) be the diagonal residual origin in
\(\operatorname{Spec}A_\theta\).  By the complete lifting property
of AC-G101 at that origin,
\(\widehat{(A_\theta)_{x_{\rm diag}}}=R_\theta\), with its
specified universal coordinates.
Every proper homogeneous prime of \(A_\theta\) is contained
in \(x_{\rm diag}\).  Indeed its degree-zero part is a proper
ideal in the local ring \(B_\theta\), hence is in
\(\mathfrak m_{B_\theta}\); all elements of nonzero degree
vanish at the origin.

Here is the required graded separation argument without assuming
that this completion is faithfully flat over the entire algebra.
Every finite graded module over a Noetherian \(\mathbb Z\)-graded
ring has a finite filtration by shifted modules \(A/P\) for
homogeneous primes \(P\).  To see this, choose a nonzero homogeneous
vector whose homogeneous annihilator is maximal.  If homogeneous
\(a,b\) satisfy \(abx=0\) and \(ax\ne0\), maximality of the
annihilator gives \(bx=0\).  Thus that annihilator is a homogeneous
prime: leading terms of finite homogeneous sums verify primality
for arbitrary elements.  Pass to the graded quotient by \(Ax\)
and repeat; the ascending-chain condition stops the construction.
All these primes lie under \(x_{\rm diag}\).
Each \(A/P\) therefore injects into its localization there, and
exactness of localization gives the same injectivity for the module.
Apply this to \(A_\theta/E_A\).
Krull separation in the Noetherian local ring then makes its
localization inject into its completion.  Exactness of completion
on finite modules identifies the latter with \(R_\theta/E_R\).
Consequently
\[
                 A_\theta/E_A\ \lhook\joinrel\longrightarrow
                                  R_\theta/E_R .
 \tag{G106.10}
\]
These two completion facts are the standard Noetherian
Artin--Rees/Krull facts already used in AC-G101; the separate
localization injection was proved above.

\emph{5. Restrict the endpoint map to the actual generators.}
Under (G106.4), every entry of the rank-two representation maps to
\(u^{-1}\) times the corresponding endpoint entry, and every
local character maps by the formula in the statement.
Thus the dense generators of \(B_2\) map into the image of
\(B_\theta\) in \(R_\theta/E_R\).
That image is closed: \(B_\theta\) is compact, the target is
Hausdorff, and the quotient map is continuous.
Taking closure places the whole of \(B_2\) in that image.
Equation (G106.10) identifies this image with
\(B_\theta/E_B\).  Since \(A_2\) is generated by \(B_2\)
and the actual matrix entries, the second map in (G106.6) follows.
Both directions preserve the matrix-entry weights.

Finally a surviving zero-weight monomial from rank three contains
only endpoint entries and middle diagonal values.  Removing its
middle factors leaves a zero-weight endpoint monomial, which is
an \(\mathcal O\)-unit multiple of one from rank two.
The local middle characters are the values of \(m\), already
proved in Step 2.  The defining closure of \(B_\theta\) therefore
proves the last assertion of the proposition.
\end{proof}

\begin{theorem}[AC-G106-even: exact finite support on all even primes]
\label{AC-G106-even}
Retain the literal rank-two input \(\mathrm{OffFin}\): there is
an ideal \(M_2\subset B_2\) such that
\[
 \mathfrak m_{B_2}\subset\sqrt{M_2+\lambda_2},\qquad
 M_2\subset P\cap B_2
   \quad\text{for every off-triangular prime }P\subset A_2,
 \tag{G106.11}
\]
where
\[
 \lambda_2=(\varpi)+
 (\lambda_{v,i}(h)-1:h\in I_v\cap\ker\gamma, v,i).
 \tag{G106.12}
\]
Equivalently one may retain exactly \(\mathrm{FinLam}(M_2)\)
and rank-two \(\mathrm{AllGood}(M_2)\), since rank-two
\(\mathrm{TI}\) means off both triangular loci.
Let \(M_E\subset B_\theta\) be the inverse image of the ideal
generated by the image of \(M_2\) in \(\overline B_\theta\).
Then
\[
 \begin{split}
 \mathfrak m_{B_\theta}&\subset
                \sqrt{M_E+\lambda_\theta},\\
 M_E&\subset P\cap B_\theta
 \quad\text{for every prime }P\subset A_\theta
       \text{ containing }E_A\text{ and off both triangular loci}.
 \end{split}
 \tag{G106.13}
\]
Here \(\lambda_\theta\) uses the joint residual inertia kernel,
exactly as in AC-G101 and AC-G102.  This is finite support on
\emph{all even primes}; it does not require their middle character
to be one.  It does not say that \(M_E\) is the particular
rank-three native automorphic ideal.
\end{theorem}

\begin{proof}
For such a prime \(P\), its inverse image under
\(A_2\to\overline A_\theta\to A_\theta/P\) is prime.
Because all odd entries vanish in this quotient, failure of both
triangular conditions means precisely that some endpoint entry
in each of the two off-diagonal positions is nonzero.
Multiplication by \(u^{-1}\) does not change this.
The pulled-back rank-two prime is therefore off triangular,
and (G106.11) gives the second line of (G106.13).

For the first line put
\(D=B_\theta/(M_E+\lambda_\theta)\).
If \(D=0\), that line is immediate; assume \(D\ne0\).
This is a complete Noetherian local ring of characteristic \(p\).
The formal block character \(m\) takes its values in
\(\overline B_\theta\), so it has values in \(D\).
On the joint kernel \(K_v=I_v\cap\ker\chi\cap\ker\gamma\),
all three local characters are one in \(D\).
The quotient \(I_v/K_v\) is finite of order prime to \(p\).
Consequently \(m|_{I_v}\), with residual value one, is trivial
by Lemma~\ref{lem:G106scalar}.
It is already unramified away from \(p\), so its image is a
quotient of the everywhere-unramified abelian pro-\(p\) class
quotient of the fixed CM field \(H\).
The exponent assertion of AC-G99, or directly global class field
theory, gives one integer \(h\ge1\) with \(m(g)^h=1\) for
all \(g\).  Writing \(h=p^a h'\) with \((h',p)=1\), the
prime-to-\(p\) factor is invertible near one and hence
\[
                       (m(g)-1)^{p^a}=0\quad\text{in }D.
 \tag{G106.14}
\]
Thus \(m=1\) in \(D_{\rm red}\), without asserting its
triviality in \(D\).

For \(h\in I_v\cap\ker\gamma\), the ratio
\(\lambda_{v,i}=u^{-1}\kappa_{v,2i}\) is trivial on \(K_v\)
and has residual value one.  Its value at \(h\) has order
dividing the order-three group
\((I_v\cap\ker\gamma)/K_v\); the scalar rigidity lemma
therefore gives \(\lambda_{v,i}(h)=1\) already in \(D\).
So the map \(B_2\to D\) kills \(M_2+\lambda_2\).
The first clause of (G106.11) says that its entire maximal ideal
has nilpotent image.  It follows that \(B_2\to D_{\rm red}\)
factors through \(k\).
By Proposition~\ref{AC-G106-corner}, \(D_{\rm red}\) is
 topologically generated by this image and the values of \(m\).
All these values lie in \(k\).  The finite set \(k\) is closed
in the Hausdorff ring \(D_{\rm red}\), so \(D_{\rm red}=k\).
This is exactly the first line of (G106.13).
\end{proof}

\begin{proposition}[AC-G106-frame: every old rational frame condition on a high layer]
\label{AC-G106-frame}
Use the original rational character \(\gamma\) on the same selected
host of AC-G93-fill, but form \(\Gamma_\gamma=G_F/N_\gamma\)
and its CM kernel as in (G106.2).  At all sufficiently high
cyclotomic layers there is an integer \(C_0\), possibly depending
on the fixed initial host and residual character, for which this
rational frame satisfies every field of the literal source
\(\mathrm{Frame.Ok}(C_0)\) and \(\mathrm{Frame.Big}(C_0)\).
The local groups remain the actual full groups.
No cubic-character \(\mu=0\), numerical uniform-host bound,
or automorphy statement is used or asserted.
\end{proposition}

\begin{proof}
\emph{The topological, local and Euler fields.}
All groups are actual profinite quotients.  Complex conjugation
has order two and CM value \(-1\); its real conjugates have the
same properties.  The CM kernel is open, the finite residual
character has open kernel, and all quotient/local/cyclotomic maps
are continuous.  Since \(\psi(-1)=(-1)^e\), one has
\(\gamma(c)=1\).  The normalization element of AC-G90 was chosen
with \(\chi(g_0)=1\) and \(\gamma(g_0)^2\ne1\), so it is also
the original rank-two normalization element.  The number of real
places is \([F:\mathbb Q]=\sum_{v\mid p}d_v\).
The arithmetic \(\Phi_p\) and Euler assertions of AC-G96 apply
with the finite residual quotient generated by \(\gamma\) and
the CM character.  Its proof permits this quotient just as it
permits the larger joint quotient: after splitting the finite
prime-to-\(p\) datum it is the same full arithmetic
\(p\)-ramification group, and descent averages over a
prime-to-\(p\) group.  In particular the Euler statement is for
all finite coefficient representations specified by the source,
not only for the two rational characters used below.
AC-G99 gives the literal \texttt{gen} and \texttt{cl} fields for
this quotient and these local maps.

The local image of \(\bar\epsilon\) on inertia has order
\(268\), unchanged by the retained \(p\)-power extensions.
The square of \(\operatorname{comp}_p\psi\) is one, so
\(\gamma^2|_{I_v}=\bar\epsilon^{2e}\ne1\), giving
\texttt{f3}.  On \(I_v\cap\ker\epsilon\), both the power of
\(\epsilon\) and the trivial or quadratic conductor-\(p\)
component of \(\psi\) are one.  This proves the full, not merely
residual-cyclotomic, field \texttt{Psi\_N}.
The characters \(\bar\epsilon\gamma^{\pm2}\eta_K\)
restrict on this inertia to \(\bar\epsilon^{1\pm2e}\).
These odd exponents are not divisible by \(268\), giving
\texttt{f2c}.  The actual scalar element of AC-G89, retained
on the selected host and each cyclotomic layer by AC-G93-fill,
has \(\gamma=1\) and \(\bar\epsilon=-1\) in the CM kernel.
It proves \texttt{f2a}, both signs of \texttt{f2b}, and
\texttt{f5} simultaneously.

\emph{Both required strict classes already exist.}
In the common prime-to-\(p\) residual splitting field used in
AC-G93-sector, consider on the real base the characters
\[
                   \gamma^{-1}\eta_K,
                    \qquad \gamma^{-2}\eta_K.
 \tag{G106.15}
\]
Both are nontrivial, since \(\gamma\) is even and \(\eta_K\)
is odd.  The quotient
\(J_k\oplus k\simeq k[\Delta]^r\), \(r\ge1\), contains every
nontrivial character sector.  Thus the very same selected extension
supplies a nonzero strict class in each of (G106.15).
No additional class-field construction is required.
The local invariant spaces vanish by \texttt{f3}, so the strict
classes are nonzero in the literal cone groups.
Restriction through a cyclotomic step is injective: a nontrivial
prime-to-\(p\) character stays nontrivial on a subgroup of
\(p\)-power index, and inflation--restriction then has zero kernel.
This supplies \texttt{f4} and \texttt{f4'} on every later layer.

\emph{All three rational upper bounds have slope zero.}
The coefficients for \texttt{b1} are the two real-field sectors
\(\gamma^{-1}\) and \(\gamma^{-1}\eta_K\), by quadratic Shapiro;
those for \texttt{b1p} are \(\gamma^{-2}\eta_K\).
They all split over a field obtained by adjoining fixed rational
abelian coefficient fields to the actual real host and \(K\).
Before the cyclic \(p\)-power host extensions, this is an abelian
extension of \(\mathbb Q\).  Ferrero--Washington gives
cyclotomic \(\mu=0\) there, as the same published input retained
in AC-G76 and AC-G98.  Each subsequent host step is cyclic of
\(p\)-power degree.  Apply AC-G100 to these rational sectors;
their downstairs \(k[[t]]\)-rank is zero, and its inequality
makes every upstairs rank zero.  A common initial cyclotomic
reindexing removes any finite intersection with the cyclotomic
tower.  Equivalently the finite-\(p\) ascent proof already given
in AC-G76 applies to these same fields.

The compact modules in AC-G74 are finite over \(k[[t]]\).
Rank zero therefore makes them finite dimensional over \(k\).
The finite-layer class-group control gives a uniform bound for
all three strict dimensions.  Its possible local permutation
corrections are uniformly bounded by the fixed number of
coefficient-prime places; after eventual total ramification that
number is fixed.  In the nontrivial rational sectors their local
invariant spaces are zero, so the cone has no additional local
\(H^0\) term.  The trivial full-local term \texttt{h1} is
uniformly bounded by the identical rational-field argument and
the exact Euler--duality formula (G76.2), as already used in
AC-G98.  Thus one fixed integer \(C\) bounds \texttt{b1},
\texttt{b1p}, and \texttt{h1}-1 on the entire later tower.

Take \(C_0=C\).  The total degree and each positive local
degree grow by \(p^n\), whereas this constant does not.
Eventually \(3C_0+7<d\) and every \(d_v>5\), exactly the two
clauses of \texttt{Big}.  The preceding paragraphs verified
all fields in the complete \texttt{Frame.Ok} declaration.
This proves the proposition with no effective bound for the
first suitable layer.
\end{proof}

\paragraph{The remaining named rank-two arithmetic input.}
The pinned theorem \path{ComponentsGoodRankTwo.lean},
\texttt{components\_good\_rank\_two}, takes precisely
\(\mathrm{Ok}\), \(\mathrm{Big}\), \(\mathrm{FinLam}(M_2)\),
\(\mathrm{P1'}(M_2)\), and \(\mathrm{SeedIn}(M_2)\), and
returns \(\mathrm{AllGood}(M_2)\).
The preceding proposition supplies its first two inputs on the
actual rational frame.  Formula (G106.11) is therefore discharged
if, and only for the purpose of this sufficient route, the remaining
three named inputs are supplied for one actual rank-two ideal.
No proof of these three inputs is hidden in (G106.4).
The rank-two native twist must have
\(\kappa_2\kappa_2^c=\epsilon^{-1}\), whereas the rank-three
\(\kappa_3=\epsilon^{-1}\) has conjugate norm \(\epsilon^{-2}\).
They cannot be substituted for one another.
The exact native twist and source arithmetic interface are a separate
supply question; no unconditional \texttt{AllGood} conclusion is
part of the present statements.  In particular,
\path{HostFrameEx.lean}, \texttt{GalLink.ab}, assumes that its CM
host is abelian over \(\mathbb Q\).  The selected class-field host
here is not asserted to have that property.  Thus the packaged
\texttt{cRealEx} and \texttt{seedPtOf} cannot simply be invoked on
this host.  The actual native \texttt{HidaSetting}, \texttt{Lev},
full local groups, and the map of AC-G97 are available; using them
directly in rank two still requires the stated finite-support and
seed interfaces to be verified.  This distinction leaves the
all-even transfer (G106.13) exact while recording its actual
arithmetic premise.

\paragraph{Source and geometric scope.}
The complete pinned source definitions used here are
\path{ContinuousCoboundaries2.lean}, \path{LeAdicClosure.lean},
\path{Integral.lean}, \path{CornerUnits.lean}, and
\path{PropagationTheoremGL2.lean}.
The complete corner proofs inspected are
\path{ExistsFrobeniusCorner.lean},
\path{FrobeniusCornerOrderSkips.lean},
\path{FrobeniusCornerMemCone.lean},
\path{FrobeniusCornerConeMap.lean}, and
\path{FrobeniusCornerMiddleTrivial.lean}.
The new proof above replaces their one-character residue calculation
by actual cubic scalar rigidity; it does not relabel the old
rank-three residual functor as the cubic one.
All source files have the pinned commit and SHA-256 receipts in the
companion manifest.  No Lean execution is claimed.

In geometric terms (G106.4) is a product of formal deformation
functors: tensor the endpoint local system by the inverse finite
cubic line, and retain the middle rank-one deformation independently.
The map on the finite-away-kernel quotients is actual descent of
local systems, not merely equality of traces.  The finite-cover
language reuses the already audited SGA~1, Expose~V, \S\S4 and~7
dictionary; no fresh archival locator or geometric automorphy theorem
is claimed.
